\section{The compressive sector of KRMHD}\label{sec:model}

Kinetic reduced magnetohydrodynamics (KRMHD) is the long-wavelength limit, $k_\perp\rho_i\to0$, of the gyrokinetic description of anisotropic fluctuations in a plasma with one ion species and isothermal electrons \citep{Schekochihin2009}. It consists of reduced magnetohydrodynamics (RMHD) for the Alfv\'enic fluctuations and a kinetic equation for the compressive ones. We write these equations in the form used by CMM \citep{Chandran2026}, whose section~3.2 and appendix~B restate and extend the analysis of S09 \citep{Schekochihin2009}. Equation numbers of CMM refer to arXiv:2607.27981v1, and those of S09 to the published article. Table~\ref{tab:conventions} compares our conventions with those of CMM, S09 and the GANDALF code \citep{Kanekar2025}, which we use for the checks of section~\ref{sec:numerics}.

\begin{table}[t]
  \centering
  \caption{Conventions of this paper compared with those of CMM \citep{Chandran2026}, S09 \citep{Schekochihin2009} and the GANDALF code \citep[version 0.6.0;][]{Kanekar2025}. Equation numbers in parentheses are those of the respective papers; ``same'' means identical to this paper.}
  \label{tab:conventions}
  \small
  \begin{tabular}{@{}p{0.27\textwidth}p{0.25\textwidth}p{0.17\textwidth}p{0.23\textwidth}@{}}
    \toprule
    \raggedright This paper & \raggedright CMM & \raggedright S09 & \raggedright GANDALF \tabularnewline
    \midrule
    \raggedright $\vth=(2T_i/m_i)^{1/2}$, $s=\vth/v_A=\beta_i^{1/2}$
      & \raggedright same ($v_{\mathrm{th}i}$)
      & \raggedright same ($v_{\mathrm{th}i}$)
      & \raggedright only $s=\beta_i^{1/2}$ enters$^{a}$ \tabularnewline[3pt]
    \raggedright $F_M(u)=\pi^{-1/2}e^{-u^2}$, $u=v_\parallel/\vth$
      & \raggedright $F_M(v_\parallel)=\pi^{-1/2}\vth^{-1}e^{-v_\parallel^2/\vth^2}$ (below B4)
      & \raggedright as CMM (below 180)
      & \raggedright same, as implied by the ladder matrix (\ref{mod:ladder}) \tabularnewline[3pt]
    \raggedright $z^\pm=\Phi\pm\Psi$; $z^+$ propagates toward $-z$
      & \raggedright Elsasser potentials not used; (3.23)--(3.24) give the same direction for $\Phi\pm\Psi$
      & \raggedright same by (21), written $\zeta^\pm$$^{b}$
      & \raggedright same$^{b}$ \tabularnewline[3pt]
    \raggedright $\Psi=-A_\parallel/(4\pi n_im_i)^{1/2}$
      & \raggedright same; $A=A_\parallel/(\rho_iB_0)=-\Psi/(\rho_iv_A)$ (3.2)
      & \raggedright same (135)
      & \raggedright $\Psi=(z^+-z^-)/2$, called $A_\parallel$ in the documentation$^{c}$ \tabularnewline[3pt]
    \raggedright $\Lambda=\Lambda^\pm$ of (\ref{eq:Lambdapm}); $\Lambda=\infty$ for $l\ge2$
      & \raggedright same (B7)
      & \raggedright same (182)
      & \raggedright input, independent of $\beta_i$ \tabularnewline[3pt]
    \raggedright $g_m=(2^mm!)^{-1/2}\int H_m(u)\,G\,\dd u$
      & \raggedright same, written $\tilde G^\pm_m$ (B27, B30)
      & \raggedright ---
      & \raggedright same$^{d}$ \tabularnewline
    \bottomrule
  \end{tabular}
  \par\medskip
  \begin{minipage}{0.96\textwidth}\footnotesize
    $^{a}$GANDALF's documentation describes $g$ as an electron distribution and its thermal speed as $(T_e/m_e)^{1/2}$. The code integrates the ion branch equation in the form (\ref{mod:hierarchy}), and its thermal-speed parameter does not enter the right-hand side.
    $^{b}$The text of S09's \S\S2.3 and~2.5 and GANDALF's documentation state the opposite direction (footnote~\ref{mod:fndir}).
    $^{c}$The documentation writes $\boldsymbol B_\perp=\zhat\times\nabla A_\parallel$, whereas the vector potential gives $\delta\boldsymbol B_\perp=-\zhat\times\nabla_\perp A_\parallel$.
    $^{d}$The documentation writes the expansion as $\sum_mg_m\psi_m(v_\parallel/\vth)$. The ladder matrix of the code is the matrix of $L$ in the basis $H_mF_M/\sqrt{2^mm!}$, so this sum is $G/F_M^{1/2}$ in the notation used here.
  \end{minipage}
\end{table}

\subsection{Alfv\'enic fields}\label{mod:alfven}

We measure velocities in units of the Alfv\'en speed, $v_A=1$, unless stated otherwise, and take the mean field $\boldsymbol B_0$ along $\zhat$. The stream function $\Phi$ and the flux function $\Psi$ give the $\boldsymbol E\times\boldsymbol B$ velocity $\uE=\zhat\times\nabla_\perp\Phi$ and the perpendicular magnetic fluctuation $\boldsymbol b_\perp\equiv\delta\boldsymbol B_\perp/(4\pi n_im_i)^{1/2}=\zhat\times\nabla_\perp\Psi$, and to first order in the fluctuation amplitude the unit vector along the magnetic field is $\bhat=\zhat+\boldsymbol b_\perp$. With the bracket $\{f,h\}=\partial_xf\,\partial_yh-\partial_yf\,\partial_xh=\zhat\cdot(\nabla f\times\nabla h)$, advection by $\uE$ is $\uE\cdot\nabla=\{\Phi,\cdot\}$ and the derivative along the field is $\gradpar\equiv\bhat\cdot\nabla=\partial_z+\{\Psi,\cdot\}$. Both $\uE$ and $\bhat$ are divergence-free. The RMHD equations are \citep[eqs~17--18]{Schekochihin2009}
\begin{equation}\label{eq:rmhd}
  \partial_t\Psi+\{\Phi,\Psi\}=\partial_z\Phi,\qquad
  \partial_t\nabla_\perp^2\Phi+\{\Phi,\nabla_\perp^2\Phi\}=\gradpar\nabla_\perp^2\Psi .
\end{equation}
The first, the induction equation $\partial_t\Psi=\gradpar\Phi$, states that the electric field along the total magnetic field vanishes; the second is the vorticity equation. Equations (\ref{eq:rmhd}) are CMM's eqs~(3.23)--(3.24), rewritten with their normalised potentials (3.2) and $\vth^2=\beta_iv_A^2$, and they do not contain the compressive fluctuations. For the Elsasser potentials $z^\pm=\Phi\pm\Psi$ (S09's $\zeta^\pm$) the linearised equations are $\partial_tz^\pm=\pm\partial_zz^\pm$, so $z^+$ propagates toward $-z$, antiparallel to $\boldsymbol B_0$, and $z^-$ propagates toward $+z$. This is the direction given by S09's eq.~(21), by CMM's eqs~(3.23)--(3.24) and by the equations that GANDALF integrates.\footnote{S09's \S2.3 writes the exact solutions with $\zeta^\mp=0$ as $\zeta^\pm=f^\pm(x,y,z\mp v_At)$, and their \S2.5 writes the slow-wave Elsasser fields in the same way. GANDALF's documentation describes $z^+$ as propagating along $+\boldsymbol B_0$. These statements give the opposite direction of propagation to the one implied by the equations. We checked the direction given here against S09's eq.~(21), against CMM's equations, and against GANDALF's right-hand side and time step.\label{mod:fndir}}

\subsection{Kinetic equation and the two branches}\label{mod:branches}

The ion thermal speed is $\vth=(2T_i/m_i)^{1/2}$, so that $s\equiv\vth/v_A=\beta_i^{1/2}$ with $\beta_i=8\pi n_iT_i/B_0^2$. We write $\tau=T_i/T_e$ and $Z=q_i/e$. We describe the compressive part of the ion distribution by the function $g$ of S09, which is CMM's $g_{09}$ (their eq.~3.26) and differs from their non-Alfv\'enic part $g_i$ by a term proportional to $\delta B/B$ (appendix~\ref{app:branches}). It obeys CMM's eq.~(3.25), the collisionless form of S09's eq.~(157),
\begin{equation}\label{mod:kinetic}
  \frac{\dd g}{\dd t}+v_\parallel\gradpar\Big[g+\Big(\frac{Z}{\tau}\frac{\delta n}{n}+\frac{v_\perp^2}{\vth^2}\frac{\delta B}{B}\Big)F_0\Big]=0,
  \qquad \frac{\dd}{\dd t}=\partial_t+\{\Phi,\cdot\},
\end{equation}
where $F_0$ is the three-dimensional Maxwellian, and $\delta n/n$ and $\delta B/B$ are the relative fluctuations of the density and of the magnetic-field strength. Quasineutrality and perpendicular pressure balance fix both fluctuations through two velocity moments of $g$ (CMM eqs~3.27--3.28; S09 eqs~158--159),
\begin{equation}\label{mod:fields}
  \overline{g}=\frac{\delta n}{n}-\frac{\delta B}{B},\qquad
  \overline{(v_\perp^2/\vth^2)\,g}=-\frac{Z}{\tau}\frac{\delta n}{n}-2\Big(1+\frac{1}{\beta_i}\Big)\frac{\delta B}{B},
  \qquad \overline{(\cdots)}\equiv\frac{1}{n_i}\int\dd^3v\,(\cdots).
\end{equation}
These relations are algebraic: $\delta n/n$ and $\delta B/B$ at a point are linear functionals of $g$ at the same point, with coefficients that do not depend on the wavenumber. At finite $k_\perp\rho_i$ the corresponding relations contain gyroaveraging operators that depend on $k_\perp\rho_i$ \citep{Schekochihin2009,Chandran2026}, and this property is lost.

Following S09 (eq.~205) and CMM (eq.~B8), we expand $g$ in Laguerre polynomials of $v_\perp^2/\vth^2$, with coefficients $G_l(v_\parallel)$. The source term in (\ref{mod:kinetic}) is linear in $v_\perp^2$, so the coefficients with $l\ge2$ obey $\dd G_l/\dd t+v_\parallel\gradpar G_l=0$ (CMM eq.~B10; S09 eqs~210 and~212, where this part of $g$ is described as a passively mixed, undamped ballistic-type mode). The coefficients with $l=0$ and $l=1$ are coupled through $\delta n/n$ and $\delta B/B$, and S09 (\S6.2.1) diagonalised this coupling. The two combinations $G^+$ and $G^-$ of $G_0$ and $G_1$ defined by CMM's eqs~(B5)--(B6) (S09 eqs~183--184; appendix~\ref{app:branches}) obey closed equations, CMM's eq.~(B4) and S09's eq.~(181),
\begin{equation}\label{mod:Gpm}
  \frac{\dd G^\pm}{\dd t}+v_\parallel\gradpar G^\pm
  =\frac{v_\parallel F_M(v_\parallel)}{\Lambda^\pm}\,\gradpar\int G^\pm(v_\parallel')\,\dd v_\parallel',
\end{equation}
where $F_M(v_\parallel)=\pi^{-1/2}\vth^{-1}e^{-v_\parallel^2/\vth^2}$ is the one-dimensional Maxwellian and
\begin{equation}\label{eq:Lambdapm}
  \Lambda^\pm=-\frac{\tau}{Z}+\frac{1}{\beta_i}\pm\kappa,\qquad
  \kappa=\Big[\Big(1+\frac{\tau}{Z}\Big)^2+\frac{1}{\beta_i^2}\Big]^{1/2}
\end{equation}
(CMM eq.~B7; S09 eqs~182 and~204).\footnote{CMM's eq.~(B16) prints $\kappa=[(1+\tau/Z)+\beta_i^{-2}]^{1/2}$, without the square on $1+\tau/Z$. With that $\kappa$ their decompositions (B15) and (B17) of the compressive free energy and phase-space helicity do not hold; with the square, as in S09's eq.~(204), they do (appendix~\ref{app:branches}). The square root in CMM's eqs~(B6)--(B7) is correct.} We call the right-hand side of (\ref{mod:Gpm}) the Landau coupling; in linear theory it produces the collisionless damping of the compressive fluctuations (S09, \S6.2.2). The coefficients $G_l$ with $l\ge2$ obey (\ref{mod:Gpm}) with $1/\Lambda=0$, and we call them the sector $\Lambda=\infty$. The compressive sector therefore consists of independent kinetic equations for $G^+$, for $G^-$ and for each $G_l$ with $l\ge2$, all advected by the same Alfv\'enic fields.

Since $\kappa>1+\tau/Z$ and $\kappa>1/\beta_i$, we have $\Lambda^+-1=\beta_i^{-1}+\kappa-(1+\tau/Z)>\beta_i^{-1}$ and $\Lambda^-=\beta_i^{-1}-\kappa-\tau/Z<-\tau/Z$, so that $\Lambda^+>1$ and $\Lambda^-<0$ for all $\beta_i>0$ and $\tau/Z>0$. S09 state both inequalities, below their eq.~(200), as the reason why the branch free energies are positive definite (section~\ref{sec:slices}). At $\beta_i=\tau=Z=1$, $\Lambda^\pm=\pm\sqrt5$. Appendix~\ref{app:branches} gives the limits of $\Lambda^\pm$ at high and low $\beta_i$.

\subsection{One branch in the variable \texorpdfstring{$u$}{u}}\label{mod:onebranch}

From here on we drop the branch label. We measure the parallel velocity by $u=v_\parallel/\vth$, write $F_M(u)=\pi^{-1/2}e^{-u^2}$ and $M_k[G]=\int u^kG\,\dd u$, and use $F_M(v_\parallel)\,\dd v_\parallel=F_M(u)\,\dd u$ and $G(v_\parallel)\,\dd v_\parallel=G(u)\,\dd u$. Since $v_\parallel=su$ when $v_A=1$, equation (\ref{mod:Gpm}) becomes
\begin{equation}\label{eq:branch}
  \partial_tG+\{\Phi,G\}+s\,\gradpar(LG)=0,
\end{equation}
with the branch operator
\begin{equation}\label{eq:L}
  LG=uG-\frac{uF_M}{\Lambda}\,M_0[G],
\end{equation}
which acts on the velocity variable only. The two branches have $\Lambda=\Lambda^\pm$, and each $G_l$ with $l\ge2$ obeys (\ref{eq:branch}) with $\Lambda=\infty$, for which $L$ is multiplication by $u$. The operator $L$ contains no wavenumber. The reason is that the field relations (\ref{mod:fields}) are algebraic, so the Landau coupling enters only through the number $\Lambda$, which is the same at every point and on every scale.

For given $\Phi$ and $\Psi$, equation (\ref{eq:branch}) is linear in $G$. In the full system its nonlinear terms are the advection $\{\Phi,G\}$ and the part $s\{\Psi,LG\}$ of the streaming, which follows the perturbed field lines. The Alfv\'enic fields obey (\ref{eq:rmhd}) whatever $G$ does, so the compressive sector is passive: it is advected by the Alfv\'enic fluctuations and does not act back on them. Passivity holds only in the limit $k_\perp\rho_i\to0$; at finite $k_\perp\rho_i$ the induction equation, for example, contains the parallel gradient of $\delta n/n$ (CMM, \S3.1).

The density and field-strength fluctuations are linear combinations of $M_0[G^+]$ and $M_0[G^-]$, given in (\ref{appA:dndB}), which are one half of the right-hand sides of S09's eqs~(202)--(203), as CMM found (footnote to their eq.~B15). The parallel ion flow involves $M_1[G^\pm]$ and the pressure perturbations involve $M_0[G^\pm]$ and $M_2[G^\pm]$ (appendix~\ref{app:branches}); the sector $\Lambda=\infty$ contributes to none of these moments.

\subsection{Hermite representation}\label{mod:hermite}

Following CMM (eqs~B27 and~B30), we expand $G=\sum_{m\ge0}g_mH_m(u)F_M(u)/\sqrt{2^mm!}$, so that $g_m=\int H_m(u)\,G\,\dd u/\sqrt{2^mm!}$, where $H_m$ are the Hermite polynomials. Equivalently, $G/F_M^{1/2}=\sum_mg_m\psi_m(u)$ with the orthonormal Hermite functions $\psi_m(u)=H_m(u)e^{-u^2/2}/(2^mm!\sqrt\pi)^{1/2}$. The coefficients $g_m$ are CMM's $\tilde G_m$ and the Hermite moments evolved by GANDALF, and $M_0[G]=g_0$, $M_1[G]=g_1/\sqrt2$. In this basis (\ref{eq:branch}) is the hierarchy
\begin{equation}\label{mod:hierarchy}
  \partial_tg+\{\Phi,g\}+s\,\gradpar(Ag)=0,\qquad g=(g_0,g_1,g_2,\dots)^{\mathsf T},
\end{equation}
where the ladder matrix $A$, the matrix of $L$, is tridiagonal with zero diagonal,
\begin{equation}\label{mod:ladder}
  A_{m,m+1}=A_{m+1,m}=\sqrt{\frac{m+1}{2}}\quad\text{except}\quad A_{10}=\frac{1-1/\Lambda}{\sqrt2}.
\end{equation}
\citet{Kanekar2015} studied the linear, one-dimensional form of this hierarchy with forcing, and \citet{Meyrand2019} integrated it in three-dimensional KRMHD simulations. The weight matrix $\Pw=\mathrm{diag}(p_m)$, with $p_0=1-1/\Lambda$ and $p_m=1$ for $m\ge1$, symmetrises it, $\Pw A=(\Pw A)^{\mathsf T}$. The quadratic form $\sum_mp_mg_m^2=\int G^2F_M^{-1}\,\dd u-M_0[G]^2/\Lambda$ is the free-energy product $\fe{G}{G}$ of section~\ref{sec:slices}; multiplied by $n_iT_i/2$ and averaged over the volume, it is CMM's branch free energy $W^\pm$ (their eqs~B13 and~B31). A truncation at $m=M$ needs a closure for $g_{M+1}$. Unless stated otherwise we use the zero closure $g_{M+1}=0$, which keeps the first $M+1$ rows and columns of $A$.

GANDALF \citep{Kanekar2025} integrates (\ref{eq:rmhd}) in Elsasser form together with (\ref{mod:hierarchy}) for one branch, with optional dissipation, hypercollisions and forcing (section~\ref{sec:lagrangian} and appendix~\ref{app:timestepping}). It truncates the hierarchy at $m=M$ with the zero closure or with $g_{M+1}=g_{M-1}$, which GANDALF calls the symmetric closure and we call the copy closure; the copy closure changes only the entry $A_{M,M-1}$ (appendix~\ref{app:truncated}). We checked to round-off that GANDALF's right-hand side with dissipation and hypercollisions switched off is (\ref{mod:hierarchy}), with the bracket and $\gradpar$ defined above, for both signs of $\Lambda$ and both closures. GANDALF takes $\Lambda$ and $\beta_i$ as separate inputs. A run therefore describes part of the compressive sector of KRMHD only if $\Lambda$ is $\Lambda^+$ or $\Lambda^-$ of (\ref{eq:Lambdapm}) for the chosen $\beta_i$ and some $\tau/Z$, for example $\Lambda=\pm\sqrt5$ at $\beta_i=\tau=Z=1$, or if $1/\Lambda=0$, which describes a Laguerre coefficient with $l\ge2$ at any $\beta_i$.

The rest of the paper concerns a single equation (\ref{eq:branch}) with $\Lambda\notin[0,1]$, which describes either branch ($\Lambda=\Lambda^\pm$) or a Laguerre coefficient with $l\ge2$ ($\Lambda=\infty$); we refer to any such equation as a branch. The compressive sector is collisionless and unforced except where collisions, hypercollisions or forcing are stated explicitly.

\begin{remark}[Formal correspondence with Vlasov--Poisson]\label{mod:vp}
For $\Phi=\Psi=0$, a solution of (\ref{eq:branch}) proportional to $e^{\ii(kz-\omega t)}$, with $k>0$ and $\mathrm{Im}\,\omega>0$, exists only if $\eps(\zeta)=1-[1+\zeta Z(\zeta)]/\Lambda=0$, with $\zeta=\omega/(ks)$; continued analytically to $\mathrm{Im}\,\omega\le0$, this is S09's dispersion relation (185). Here $\eps$ is the dielectric function (\ref{eq:eps}) of section~\ref{sec:slices}, and $Z(\zeta)=\int F_M(x)\,(x-\zeta)^{-1}\dd x$ is the plasma dispersion function of \citet{FriedConte1961}, not the charge number. For a Maxwellian with thermal speed $\vth=(2T/m)^{1/2}$ and with $\zeta=\omega/(k\vth)$, the dielectric function of the electrostatic Vlasov--Poisson system is $1+[1+\zeta Z(\zeta)]/(k\lambda_D)^2$, with $\lambda_D$ the Debye length, and that of a self-gravitating collisionless system is $1-[1+\zeta Z(\zeta)]/(k/k_J)^2$, with $k_J$ the Jeans wavenumber. A branch with $\Lambda<0$ therefore has the Vlasov--Poisson dielectric function at the fixed value $(k\lambda_D)^2=|\Lambda|$, and a branch with $\Lambda>1$ has that of a self-gravitating system on the stable side of the Jeans threshold, at $(k/k_J)^2=\Lambda$. The range $0<\Lambda<1$, which no KRMHD branch reaches, corresponds to the Jeans-unstable side, and $\Lambda=1$ to the marginal case. The KRMHD branches differ from both systems in that $\Lambda$ does not depend on $k$, which is the property used in section~\ref{sec:slices}.
\end{remark}
